\begin{frame}
\frametitle{Cíl přednášky}

Porozumět implementaci jednoduchého počítače, který se skládá z procesoru a oddělené paměti pro instrukce a data.

Naším cílem je implementovat následující instrukce:

\begin{itemize}
\item Číst a zapisovat hodnotu z/do datové paměti.
\begin{itemize}
\item \textbf{\texttt{lw}} - načtení slova, \textbf{\texttt{sw}} - uložení slova
\end{itemize}
\item aritmetické a logické instrukce: \textbf{\texttt{add}}, \textbf{\texttt{sub}}, \textbf{\texttt{and}}, \textbf{\texttt{or}}, \textbf{\texttt{slt}}
\begin{itemize}
\item Immediátové varianty: \textbf{\texttt{addi}}, \textbf{\texttt{ori}}, horní bity načítají \textbf{\texttt{olui}},\textbf{\texttt{auipc}}
\end{itemize}
\item Instrukce pro změnu toku programu/skok \textbf{\texttt{beq}}
\item Volání podprogramu \textbf{\texttt{jal}}, \textbf{\texttt{jalr}} (zajišťuje i návrat z podprogramu \textbf{\texttt{jr~ra}})
\item CPU se bude skládat z řídicí jednotky a ALU.
\item Poznámky:
\begin{itemize}
\item Implementace bude minimální (jednocyklový procesor - všechny operace budou probíhat v jednom cyklu, tzn. zpracovány v jednom kroku/hodině)
\item Přednáška 5 je zaměřena na efektivnější realističtější implementaci zřetězenného (pipelined) CPU.
\end{itemize}
\end{itemize}

\end{frame}
